% Spivak Capítulo 4, Problema 5
\section*{Spivak Capítulo 4, Problema 5}
\addcontentsline{toc}{section}{Spivak Capítulo 4, Problema 5}

Dibujar los conjuntos de puntos $(x,y)$ que satisfacen cada una de las
relaciones indicadas. La sugerencia de Spivak es intercambiar $x$ e $y$;
esto nos permite reconocer inmediatamente curvas ya vistas en el capítulo
anterior.

\begin{enumerate}[label=(\roman*)]
\item $x=y^{2}$.

  Aquí $x$ está expresado en función de $y$; para cada valor de $y$ se
  obtiene un único $x$ no negativo. La curva es la parábola clásica con
  eje horizontal y vértice en el origen. Si intercambiamos variables
  obtenemos $y=x^{2}$, que es la parábola con eje vertical conocida de
  antes; así confirmamos que se abre hacia la derecha.

  % dibujo
  \begin{center}
  \begin{tikzpicture}[scale=0.7]
    \draw[->] (-1,0) -- (4,0) node[right]{$x$};
    \draw[->] (0,-2) -- (0,2) node[above]{$y$};
    \draw[domain=-2:2,smooth,variable=\t,blue] plot ({\t*\t},{\t});
  \end{tikzpicture}
  \end{center}

\item $\displaystyle\frac{y^{2}}{a^{2}}-\frac{x^{2}}{b^{2}}=1$.

  Esta ecuación se reconoce como la hipérbola estándar con ejes alineados
  con los ejes coordenados, pero con los papeles de $x$ e $y$ cambiados en
  comparación con la forma $x^{2}/a^{2}-y^{2}/b^{2}=1$ estudiada en el
  capítulo anterior. Por lo tanto la hipérbola tiene ramas que se abren
  hacia arriba y hacia abajo; el parámetro $a$ controla la distancia del
  vértice al origen sobre el eje $y$, y $b$ la asíntota se aproxima al
  eje horizontal. En el plano la curva es la imagen rotada 90° de la
  hipérbola "horizontal".

  % dibujo con a=2,b=1
  \begin{center}
  \begin{tikzpicture}[scale=0.7]
    \draw[->] (-2,0) -- (2,0) node[right]{$x$};
    \draw[->] (0,-3) -- (0,3) node[above]{$y$};
    \draw[domain=1.4:3,smooth,variable=\t,blue] plot ({\t/2},{\t});
    \draw[domain=-3:-1.4,smooth,variable=\t,blue] plot ({\t/2},{\t});
  \end{tikzpicture}
  \end{center}

\item $x=|y|$.

  La ecuación se descompone en dos semirrectas: cuando $y\ge0$ da $x=y$,
  y cuando $y\le0$ da $x=-y$. Geométricamente es el doble cono con vértice
  en $(0,0)$, abriéndose hacia la derecha; intercambiando $x$ e $y$ se
  obtiene el mismo dibujo porque $|x|=y$ representa el cono hacia arriba y
  abajo.

  \begin{center}
  \begin{tikzpicture}[scale=0.7]
    \draw[->] (-1,0) -- (3,0) node[right]{$x$};
    \draw[->] (0,-3) -- (0,3) node[above]{$y$};
    \draw[blue] (0,0) -- (2,2);
    \draw[blue] (0,0) -- (2,-2);
  \end{tikzpicture}
  \end{center}

\item $x=\sin y$.

  Para cada $y$ real el valor de $x$ se sitúa entre $-1$ y $1$; la gráfica
  es la famosa curva ondulada que cruza el eje $x$ con periodo $2\pi$.
  Si se quiere ver con más claridad, se puede despejar $y=\arcsin x$, en
  cuyo caso aparece la función inversa estudiada anteriormente, excepto
  que ahora $y$ no está restringido al intervalo $[-\pi/2,\pi/2]$ y existe
  una infinidad de $y$ para cada $x\in[-1,1]$. Alternativamente, pensar en
  la sustitución $u=y$ de la parábola $u=\sin x$ estudiada antes da la
  misma forma tras intercambiar los ejes.

  \begin{center}
  \begin{tikzpicture}[scale=0.7]
    \draw[->] (-1.5,0) -- (1.5,0) node[right]{$x$};
    \draw[->] (0,-3.5) -- (0,3.5) node[above]{$y$};
    \draw[domain=-3.14:3.14,smooth,variable=\t,blue] plot ({sin(\t)},{\t});
  \end{tikzpicture}
  \end{center}
\end{enumerate}

Estas descripciones son suficientes para trazar las curvas con lápiz y
papel; la indicación de intercambiar variables facilita el reconocimiento.  